%%% - - -   UNCLASSIFIED   - - - %%%

\documentclass[11pt, letterpaper]{article}

\newcommand{\thisdocument}{MILSTD}

\usepackage{format/NGICStyle}
\usepackage{format/DurandalBib}
\usepackage{format/NGICGlossary}
\usepackage{values/T150ModelValues}

\graphicspath{{images}}
\newcommand\csvpath{csv}

\makeatletter
\def\input@path{{content}}
\makeatother

\chead{\documentlabel}
\cfoot{\documentlabel\ \\ \thepage}
\presetkeys{todonotes}{disable}{}

\begin{document}

\pagenumbering{roman}

\InDevelopmentfalse

\ifInDevelopment
	\draft
	\linenumbers
	\colorlet{ModelValueColor}{blue!50!cyan}
\fi

%% ============================================================
%%  FRONT MATTER
%% ============================================================

\begin{center}
    \vspace*{0.3in}
    \Huge\textbf{\LARGE\unclass Accreditation Report\\}
    \vspace*{0.15in}
    \Large\textbf{\unclass T-150/TRV-150 Heavy-Lift UAS\\AFSIM Simulation Model v2.0\\}
    \normalsize\textbf{\unclass AFSIM Project 005\\}
    \vspace*{0.3in}
    \includegraphics[width=4in]{images/ngic}
    \vspace*{0.1in}
    \Large
    \textbf{Version: 1.0 \\
    Date: March 5, 2026 \\
    Classification: UNCLASSIFIED \\}
\end{center}

\vfill

\tableTwoNow{\unclasslabel}{\unclass Document Information}{accrpt-docinfo}{l,2.5in}{l,3.5in}
{
    Field & Value \\
    Document Version & 1.0 \\
    Date & March 5, 2026 \\
    Classification & UNCLASSIFIED \\
    Sponsoring Organization & [PLACEHOLDER - Organization Name] \\
    Accreditation Authority & [PLACEHOLDER - Name/Title] \\
    Document Preparer & [PLACEHOLDER - Organization/Contract] \\
    Distribution Statement & [PLACEHOLDER - Distribution Statement] \\
    M\&S Name and Version & T-150 UAS AFSIM Plugin Suite v2.0 \\
    Accreditation Decision & CONDITIONAL ACCREDITATION \\
}

\unclass This document is prepared in accordance with:
\begin{itemize}
    \item \unclass \textbf{DoDI 5000.61}, ``DoD Modeling and Simulation Verification, Validation, and Accreditation,'' September 17, 2024 [1]
    \item \unclass \textbf{MIL-STD-3022}, Appendix D, ``Documentation of VV\&A for Models and Simulations,'' January 28, 2008 [2]
    \item \unclass \textbf{NGIC 4.0 Documentation Templates}, January 30, 2026
\end{itemize}

\pagebreak

\begin{center}
    \begin{tblr}
        {
            colspec     ={|Q[c,0.8in]|Q[c,0.8in]|Q[4in]|},
            hline{1-Z}  ={.5pt},
            cell{1}{1}  ={c=3}{halign=l},
            cell{Z}{1}  ={c=3}{halign=r},
            row{2}      ={bg=TableBlue,fg=white,halign=c},
        }
        \unclasslabel\\
        Version & Date     & Changes \\
        1.0     & 2026-03-05 & Initial release. Accreditation assessment for T-150 UAS AFSIM Plugin Suite v2.0 covering all six development phases. \\
        \unclasslabel\\
    \end{tblr}
\end{center}

\vfill\pagebreak

\tableofcontents\thispagestyle{fancy}\pagebreak
\listoftables\thispagestyle{fancy}\pagebreak

\pagenumbering{arabic}
\cfoot{\classlabel\ \\ {Page \thepage\ of \pageref{LastPage}}}


%% ============================================================
%%  EXECUTIVE SUMMARY
%% ============================================================

\section{\unclass Executive Summary}

\unclass This Accreditation Report documents the accreditation assessment and decision for the T-150/TRV-150 Heavy-Lift UAS AFSIM Simulation Model v2.0 (hereafter ``the Model''). The Model comprises five custom C++ plugins totaling approximately 7,736 lines of code across 26 source files, developed for AFSIM version 2.2505. The plugins collectively model the flight dynamics, battery energy system, propulsion, Vortex Ring State (VRS), and slung load dynamics of the T-150 coaxial octocopter UAS.

\unclass The accreditation assessment evaluated the Model against 13 requirements derived from the intended use of intelligence analysis of T-150 UAS operational capabilities. Verification and validation testing demonstrated that 12 of 13 requirements are fully met, with one requirement (R-009, Flight Endurance) partially met due to a calibration gap between simulated endurance (14.3~minutes) and manufacturer-specified endurance (36~minutes).

\unclass \textbf{Overall Recommendation: CONDITIONAL ACCREDITATION} for intelligence analysis use, subject to the following key conditions:

\begin{itemize}
    \item \unclass Analysts must apply an endurance scaling factor when interpreting time-based results
    \item \unclass Results should be used for comparative analysis (relative differences) rather than absolute performance predictions
    \item \unclass VRS and slung load results should be treated as indicative of behavioral trends, not as predictive outputs
    \item \unclass The Model should be recalibrated if T-150 flight test data becomes available
\end{itemize}

\unclass The Model demonstrates physics-based fidelity sufficient for intelligence-grade analysis of T-150 operational capabilities, vulnerability assessment, and counter-UAS scenario development. Energy conservation has been validated to less than 0.01\% error, subsystem interactions are correctly modeled, and behavioral trends align with established aerodynamic theory. The primary limitation is the absence of flight test data for direct validation, which is typical for threat system modeling based on open-source intelligence.

\pagebreak

%% ============================================================
%%  SECTION 1 -- Problem Statement
%% ============================================================

\section{\unclass Problem Statement}

\unclass The intelligence community requires a physics-based simulation of the T-150/TRV-150 Heavy-Lift UAS to answer analytical questions regarding the operational capabilities, performance envelope, and limitations of this heavy-lift multirotor platform. These questions support counter-UAS assessment, vulnerability analysis, and intelligence production.

\unclass Native AFSIM kinematic movers lack the fidelity required to model rotor aerodynamics, battery depletion, motor thermal dynamics, Vortex Ring State, ground effect, and slung load interactions. A custom plugin suite was developed to provide parametric-level fidelity for these physics domains.

\unclass The intended use of the simulation is intelligence analysis, specifically:

\begin{itemize}
    \item \unclass Assessing T-150 operational capabilities and performance boundaries
    \item \unclass Comparative flight pattern efficiency analysis
    \item \unclass Vulnerability assessment (VRS entry conditions, battery depletion profiles, thermal limits)
    \item \unclass Payload impact analysis on flight performance
    \item \unclass Counter-UAS (C-UAS) scenario development and evaluation
\end{itemize}

\pagebreak

%% ============================================================
%%  SECTION 2 -- Requirements and Acceptability Criteria
%% ============================================================

\section{\unclass M\&S Requirements and Acceptability Criteria}

\subsection{\unclass Requirements Traceability with Final Assessment}

\tableFivNow{\unclasslabel}{\unclass Requirements Final Assessment}{accrpt-reqs}{c,0.4in}{l,1.2in}{l,1.2in}{c,0.5in}{l,2in}
{
    Req ID & Description & Acceptability Criteria & Met? & Assessment Summary \\
    R-001 & Battery SOC accuracy & SOC tracking error $<$ 1\% & \textbf{YES} & Energy conservation $<$ 0.01\% error in Test 09 \\
    R-002 & Voltage sag model & Terminal voltage follows Li-Ion curve & \textbf{YES} & SOC-dependent resistance produces correct sag (Test 06) \\
    R-003 & Motor thermal dynamics & Stable equilibrium; no runaway & \textbf{YES} & Stable equilibrium under sustained load (Test 05) \\
    R-004 & Position hold stability & Drift $<$ 0.5~m/s in hover & \textbf{YES} & PID controller maintains hover $\approx$ 0~m/s drift (Test 01) \\
    R-005 & Velocity tracking & Error $<$ 1~m/s during cruise & \textbf{YES} & Achieves commanded velocities within spec (Test 02) \\
    R-006 & Numerical stability & No NaN, Inf, or divergence & \textbf{YES} & No instabilities under extreme inputs (Test 08) \\
    R-007 & VRS detection & Correct state transitions & \textbf{YES} & Correct transitions based on descent rate and airspeed (Test 03) \\
    R-008 & Ground effect & Thrust augmentation near ground & \textbf{YES} & Cheeseman-Bennett model produces correct increase (Test 07) \\
    R-009 & Flight endurance & Within 20\% of manufacturer spec & \textbf{PARTIAL} & 14.3~min vs.\ 36~min; behavioral trends correct; scaling required \\
    R-010 & Slung load dynamics & Correct pendulum frequency/damping & \textbf{YES} & Frequency consistent with cable length parameters (Test 10) \\
    R-011 & Attitude recovery & Recovery within 5~s from disturbance & \textbf{YES} & PID recovers from disturbances within specification (Test 04) \\
    R-012 & Cruise speed & Achieves 30~m/s & \textbf{YES} & 30~m/s achieved and sustained (Test 02) \\
    R-013 & Max payload capacity & Correctly limits to 68~kg & \textbf{YES} & 68~kg modeled; performance degradation consistent (Test 11) \\
}

\subsection{\unclass Requirements Assessment Summary}

\begin{itemize}
    \item \unclass \textbf{Requirements Fully Met}: 12 of 13 (92.3\%)
    \item \unclass \textbf{Requirements Partially Met}: 1 of 13 (7.7\%) -- R-009 (Flight Endurance)
    \item \unclass \textbf{Requirements Not Met}: 0 of 13 (0.0\%)
\end{itemize}

\unclass The single partially-met requirement (R-009) is attributable to a calibration gap in the power consumption model rather than a fundamental modeling error. The behavioral trend (monotonic SOC decline with correct voltage sag) is correct; the absolute rate of energy consumption is higher than the real system, likely due to conservative efficiency assumptions. This gap is well-documented and can be addressed through analyst-applied scaling factors.

\pagebreak

%% ============================================================
%%  SECTION 3 -- Assumptions, Capabilities, Limitations, Risks
%% ============================================================

\section{\unclass M\&S Assumptions, Capabilities, Limitations, and Risks/Impacts}

\subsection{\unclass Key Assumptions}

\begin{enumerate}
    \item \unclass \textbf{Rigid body assumption}: The airframe is modeled as a rigid body. Structural flex, vibration modes, and aeroelastic effects are not represented.
    \item \unclass \textbf{Quasi-steady aerodynamics}: Rotor aerodynamics use momentum theory with quasi-steady assumptions.
    \item \unclass \textbf{Uniform inflow}: Each rotor is modeled with uniform inflow. BEMT radial variation is not implemented.
    \item \unclass \textbf{Estimated inertia tensor}: Moments of inertia are estimated from geometric assumptions and mass distribution.
    \item \unclass \textbf{Standard atmosphere}: Air density varies with altitude per ISA model.
    \item \unclass \textbf{Idealized battery cells}: All battery cells are assumed identical.
    \item \unclass \textbf{No wind turbulence}: The current model does not include atmospheric turbulence or gust modeling.
    \item \unclass \textbf{Open-source parameters}: All model parameters are derived from OSINT.
\end{enumerate}

\subsection{\unclass Model Capabilities}

\begin{itemize}
    \item \unclass \textbf{6DOF flight dynamics}: Full six-degree-of-freedom rigid body physics with Newton-Euler integration, PID attitude and velocity control, and waypoint navigation
    \item \unclass \textbf{Battery energy modeling}: Coulomb counting SOC tracking, SOC-dependent OCV, load-dependent internal resistance and voltage sag, low-battery warnings
    \item \unclass \textbf{Propulsion modeling}: Per-motor RPM and thrust calculation, motor thermal dynamics with convective cooling, dynamic power from flight state
    \item \unclass \textbf{VRS detection}: State machine detecting VRS entry conditions with configurable thresholds
    \item \unclass \textbf{Ground effect}: Cheeseman-Bennett thrust augmentation model
    \item \unclass \textbf{Slung load dynamics}: Spherical pendulum model with force feedback through AuxData
    \item \unclass \textbf{Inter-plugin communication}: Loosely-coupled blackboard architecture via AuxData
\end{itemize}

\subsection{\unclass Model Limitations}

\begin{enumerate}
    \item \unclass \textbf{Endurance calibration gap}: Simulated endurance (14.3~min) is $\approx$40\% of manufacturer specification (36~min). All time-based results require scaling.
    \item \unclass \textbf{No flight test validation}: No T-150 flight test data available for direct comparison.
    \item \unclass \textbf{Developer self-testing}: V\&V performed by development team; no independent V\&V agent.
    \item \unclass \textbf{Estimated parameters}: Several parameters estimated from class-level data.
    \item \unclass \textbf{Informational subsystems}: VRS and slung load are informational -- they report state but do not dynamically alter flight trajectory.
    \item \unclass \textbf{No degraded-mode modeling}: Motor failure, partial power loss, GPS denial not modeled.
    \item \unclass \textbf{Single environmental condition}: ISA standard conditions only.
\end{enumerate}

\subsection{\unclass Accreditation-Specific Risk Assessment}

\tableThrNow{\unclasslabel}{\unclass Accreditation Risk Assessment}{accrpt-riskassess}{l,2.5in}{c,1in}{l,2.5in}
{
    Risk Category & Risk Level & Justification \\
    Relative performance comparisons & \textbf{LOW} & Physics-based trends correct; comparative analysis well-supported \\
    Absolute performance predictions & \textbf{MEDIUM} & Endurance gap introduces quantitative uncertainty \\
    Time-critical mission analysis & \textbf{HIGH} & Without scaling, time-based predictions underestimate mission duration \\
    Behavioral trend analysis & \textbf{LOW} & Physics-based models produce correct qualitative behavior \\
    C-UAS engagement analysis & \textbf{LOW-MED} & Flight dynamics sufficient for geometry analysis; signatures not modeled \\
}

\pagebreak

%% ============================================================
%%  SECTION 4 -- Accreditation Assessment
%% ============================================================

\section{\unclass Accreditation Assessment}

\subsection{\unclass Assessment Summary}

\unclass The accreditation assessment evaluated the T-150 UAS AFSIM Plugin Suite v2.0 against 13 model requirements derived from the intended use of intelligence analysis. The assessment considered V\&V test results, documentation completeness, configuration management practices, and overall suitability.

\unclass \textbf{Result}: 12 of 13 requirements fully met. One requirement (R-009) partially met with documented calibration gap. The Model demonstrates sufficient fidelity for intelligence analysis when used with appropriate analyst awareness of its limitations.

\subsection{\unclass Strengths}

\begin{enumerate}
    \item \unclass \textbf{Physics-based modeling provides correct behavioral trends}: Momentum theory, Newton-Euler dynamics, Coulomb counting, and first-principles thermal modeling ensure physically consistent results. Energy conservation validated to $<$0.01\% error.
    \item \unclass \textbf{Energy conservation validated to high precision}: Test 09 demonstrated $<$0.01\% energy balance error over full discharge cycles.
    \item \unclass \textbf{Multiple subsystem interactions correctly modeled}: The loosely-coupled plugin architecture successfully models interactions between all five subsystems.
    \item \unclass \textbf{Configuration management via git provides traceability}: All source code, scenarios, and documentation maintained with full commit history.
    \item \unclass \textbf{Comprehensive documentation suite}: ADD v3.0, UAM v3.0, ICD v1.0, and TDD provide full traceability.
    \item \unclass \textbf{Modular architecture enables incremental confidence}: Each plugin independently verifiable. Six-phase development allowed progressive confidence building.
    \item \unclass \textbf{Boundary condition testing demonstrates robustness}: No numerical instabilities under extreme conditions.
\end{enumerate}

\subsection{\unclass Weaknesses}

\begin{enumerate}
    \item \unclass \textbf{No flight test data for direct validation}: Validation relies on specifications, literature, and physics-based consistency. Common for threat system modeling.
    \item \unclass \textbf{Endurance gap requires analyst awareness}: 2.5x discrepancy means time-based outputs cannot be used at face value.
    \item \unclass \textbf{Developer self-testing}: Risk of confirmation bias in test design and interpretation.
    \item \unclass \textbf{Some estimated parameters}: Inertia, thermal mass, and resistance coefficients estimated from class-level data.
    \item \unclass \textbf{VRS and slung load are informational only}: Report state but do not impose full aerodynamic coupling.
\end{enumerate}

\subsection{\unclass Risk Assessment}

\unclass Overall risk for intended use of intelligence analysis: \textbf{LOW to MEDIUM}.

\begin{itemize}
    \item \unclass \textbf{LOW risk} for comparative analysis and behavioral trend identification (primary intended use)
    \item \unclass \textbf{MEDIUM risk} for quantitative performance predictions (require scaling factors)
    \item \unclass \textbf{HIGH risk} for time-critical predictions without endurance scaling (mitigated by documented conditions)
\end{itemize}

\pagebreak

%% ============================================================
%%  SECTION 5 -- Accreditation Recommendations
%% ============================================================

\section{\unclass Accreditation Recommendations}

\subsection{\unclass Recommendation}

\unclass \textbf{CONDITIONAL ACCREDITATION}

\unclass The T-150/TRV-150 Heavy-Lift UAS AFSIM Simulation Model v2.0 is recommended for conditional accreditation for use in intelligence analysis of T-150 UAS operational capabilities. The conditions in Section 5.2 must be satisfied for all analytical uses.

\subsection{\unclass Conditions of Accreditation}

\unclass \textbf{(a) Endurance Scaling Requirement}: Analysts must apply an endurance scaling factor of approximately 2.5x when interpreting time-based simulation results. Simulated flight durations should be multiplied by 2.5 to estimate real-world endurance. This factor must be documented in all analytical products referencing simulation-derived time or range values.

\unclass \textbf{(b) Comparative Analysis Preference}: Results should be used for comparative analysis (relative differences between configurations) rather than absolute performance predictions. The Model's strength lies in correctly capturing the direction and relative magnitude of performance changes.

\unclass \textbf{(c) VRS and Slung Load Interpretation}: VRS detection and slung load dynamics results should be treated as indicative of behavioral trends, not as predictive of specific outcomes. The VRS processor identifies conditions consistent with VRS entry but does not impose full aerodynamic coupling.

\unclass \textbf{(d) Recalibration on Data Availability}: The Model should be recalibrated if T-150 flight test data, detailed motor efficiency curves, or measured battery discharge profiles become available. Priority calibration targets: motor efficiency vs.\ RPM curves, battery internal resistance vs.\ SOC and temperature, and propeller thrust coefficient vs.\ advance ratio.

\subsection{\unclass Approved Uses}

\unclass The Model is approved for the following analytical applications, subject to the conditions in Section 5.2:

\begin{itemize}
    \item \unclass \textbf{Intelligence analysis} of T-150 UAS operational capabilities, performance envelope, and system limitations
    \item \unclass \textbf{Comparative flight pattern efficiency analysis}, including mission profiles, waypoint optimization, and altitude/speed trade-offs
    \item \unclass \textbf{Vulnerability assessment}, including VRS susceptibility, battery depletion profiles, motor thermal limits, and operational boundary identification
    \item \unclass \textbf{Payload impact analysis}, evaluating payload mass effects on performance, endurance, and maneuverability
    \item \unclass \textbf{Counter-UAS (C-UAS) scenario development}, providing representative UAS flight dynamics for detection, tracking, and engagement analysis
    \item \unclass \textbf{Concept of Operations (CONOPS) evaluation}, assessing tactical employment concepts for heavy-lift UAS operations
    \item \unclass \textbf{Training and familiarization}, demonstrating heavy-lift UAS flight characteristics
\end{itemize}

\subsection{\unclass Prohibited Uses}

\unclass The Model is \textbf{not approved} for:

\begin{itemize}
    \item \unclass \textbf{Real-time flight control} or any application involving actual aircraft operations
    \item \unclass \textbf{Safety-critical decision making} where model inaccuracy could result in loss of life or equipment
    \item \unclass \textbf{Acquisition decisions} based solely on model output without independent data sources
    \item \unclass \textbf{Absolute performance guarantees} without application of the endurance scaling factor and documented uncertainty bounds
    \item \unclass \textbf{Prediction of specific failure modes} beyond the behavioral indicators modeled
    \item \unclass \textbf{Certification or airworthiness assessment} of any UAS system
\end{itemize}

\pagebreak

%% ============================================================
%%  SECTION 6 -- Key Participants
%% ============================================================

\section{\unclass Key Participants}

\tableTwoNow{\unclasslabel}{\unclass Key Participants}{accrpt-participants}{l,2.5in}{l,3.5in}
{
    Role & Name/Organization \\
    Accreditation Authority & [PLACEHOLDER - Name, Title, Organization] \\
    M\&S Proponent & [PLACEHOLDER - Name, Title, Organization] \\
    V\&V Agent & [PLACEHOLDER - Name, Title, Organization] \\
    SME (UAS) & [PLACEHOLDER - Name, Title, Organization] \\
    SME (Aerodynamics) & [PLACEHOLDER - Name, Title, Organization] \\
    Model Developer & [PLACEHOLDER - Name, Title, Organization] \\
    Configuration Manager & [PLACEHOLDER - Name, Title, Organization] \\
    Analyst Representative & [PLACEHOLDER - Name, Title, Organization] \\
}

\pagebreak

%% ============================================================
%%  SECTION 7 -- Actual Resources Expended
%% ============================================================

\section{\unclass Actual Accreditation Resources Expended}

\tableFouNow{\unclasslabel}{\unclass Resources Expended}{accrpt-resources}{l,2in}{c,1in}{c,1in}{l,2in}
{
    Resource Category & Planned & Actual & Notes \\
    Calendar Duration & [PLACEHOLDER] & $\sim$57 days & Development start to accreditation package \\
    Active Development Hours & [PLACEHOLDER] & $\sim$20 hours & AI-assisted development reduced effort \\
    V\&V Testing Hours & [PLACEHOLDER] & [PLACEHOLDER] & Includes scenario development, execution, analysis \\
    Documentation Hours & [PLACEHOLDER] & [PLACEHOLDER] & AI-assisted documentation reduced effort \\
    Review and Assessment Hours & [PLACEHOLDER] & [PLACEHOLDER] & SME review, accreditation meetings \\
    Total Labor (person-hours) & [PLACEHOLDER] & [PLACEHOLDER] & Sum of all categories \\
}

\unclass \textbf{Note}: AI-assisted development using Claude Code CLI substantially reduced both code development and documentation timelines. Refer to the Impact Statement [12] for detailed discussion.

\pagebreak

%% ============================================================
%%  SECTION 8 -- Lessons Learned
%% ============================================================

\section{\unclass Accreditation Lessons Learned}

\subsection{\unclass AI-Assisted Development}

\unclass AI-assisted development significantly accelerated model creation, producing over 7,700 lines of C++ code and comprehensive VV\&A documentation within approximately 20 hours. However, this approach required human calibration oversight -- the AI-generated code was architecturally sound but required manual intervention for motor efficiency calibration.

\unclass \textbf{Recommendation}: AI-assisted M\&S development should include explicit calibration review gates where SMEs validate key performance parameters before integration testing.

\subsection{\unclass AFSIM Plugin Architecture}

\unclass The AFSIM plugin architecture enabled modular V\&V. The AuxData blackboard communication pattern simplified interface testing. This modularity directly supported the six-phase iterative development approach.

\unclass \textbf{Recommendation}: Future AFSIM-based M\&S efforts should adopt modular plugin architecture with AuxData communication to facilitate incremental V\&V.

\subsection{\unclass Open-Source Intelligence as Basis for Threat Modeling}

\unclass OSINT provided sufficient data for parametric-level fidelity. The AAIB report, patent filing, motor specs, and product data collectively provided approximately 65 system parameters.

\unclass \textbf{Recommendation}: OSINT-based threat models should document source and confidence level for every parameter.

\subsection{\unclass Physics-Based Validation}

\unclass Physics-based models provide inherent validation through energy conservation and dimensional analysis. The $<$0.01\% energy accounting accuracy provides confidence in the underlying physics engine even without flight test data.

\unclass \textbf{Recommendation}: For threat system models where flight test data is unavailable, physics-based consistency checks should be a primary validation technique.

\subsection{\unclass Iterative Development and Progressive Confidence}

\unclass The six-phase approach allowed progressive confidence building. Each phase added capability while regression testing preserved previously validated behavior.

\unclass \textbf{Recommendation}: Complex M\&S efforts should adopt iterative development with defined phase gates.

\subsection{\unclass Configuration Management}

\unclass Git and GitHub provided excellent traceability. Every code change is linked to a specific commit, branch, and pull request. The development history is fully auditable.

\unclass \textbf{Recommendation}: All M\&S efforts subject to VV\&A should use version control with branch-based workflows.

\pagebreak

%% ============================================================
%%  APPENDICES
%% ============================================================

\appendix

\section{\unclass M\&S Description}

\tableFivNow{\unclasslabel}{\unclass M\&S Components}{accrpt-components}{l,1.2in}{l,1.2in}{l,1.5in}{c,0.5in}{l,1.5in}
{
    Plugin & Library & Registered Type & LOC & Purpose \\
    Rotorcraft 6DOF Mover & rotorcraft\_6dof & WSF\_ROTORCRAFT\_6DOF\_MOVER & 2,939 & Physics-based 6DOF flight dynamics with PID control \\
    Battery Processor & t150\_battery & T150\_BATTERY\_PROCESSOR & 1,139 & Li-Ion battery SOC tracking and voltage sag \\
    Propulsion Processor & t150\_propulsion & T150\_PROPULSION\_PROCESSOR & 1,281 & Motor power, thermal dynamics, RPM, ground effect \\
    VRS Processor & t150\_vrs & T150\_VRS\_PROCESSOR & 697 & VRS detection and state machine \\
    Slung Load Processor & t150\_slung\_load & T150\_SLUNG\_LOAD\_PROCESSOR & 825 & Cable-suspended payload pendulum dynamics \\
    \textbf{Total} & & & \textbf{7,736} & \\
}

\unclass The plugins communicate through the AFSIM AuxData key-value store. For full architecture description, refer to the ADD v3.0 [7].

\pagebreak

\section{\unclass Requirements Traceability Matrix (Final)}

\tableFivNow{\unclasslabel}{\unclass Final Requirements Traceability}{accrpt-rtm}{c,0.5in}{l,1in}{c,0.7in}{l,1.5in}{c,1in}
{
    Req ID & V\&V Plan Ref & Test ID & Result & Status \\
    R-001 & Sec 3.1 & Test 09 & \PASS{} ($<$ 0.01\% error) & ACCEPTED \\
    R-002 & Sec 3.2 & Test 06 & \PASS & ACCEPTED \\
    R-003 & Sec 3.3 & Test 05 & \PASS & ACCEPTED \\
    R-004 & Sec 3.4 & Test 01 & \PASS{} (drift $\sim$ 0~m/s) & ACCEPTED \\
    R-005 & Sec 3.5 & Test 02 & \PASS{} (error $<$ 1~m/s) & ACCEPTED \\
    R-006 & Sec 3.6 & Test 08 & \PASS{} (no NaN) & ACCEPTED \\
    R-007 & Sec 3.7 & Test 03 & \PASS & ACCEPTED \\
    R-008 & Sec 3.8 & Test 07 & \PASS & ACCEPTED \\
    R-009 & Sec 3.9 & Test 12 & PARTIAL (14.3 vs 36~min) & ACCEPTED WITH CONDITIONS \\
    R-010 & Sec 3.10 & Test 10 & \PASS & ACCEPTED \\
    R-011 & Sec 3.11 & Test 04 & \PASS & ACCEPTED \\
    R-012 & Sec 3.12 & Test 02 & \PASS{} (30~m/s achieved) & ACCEPTED \\
    R-013 & Sec 3.13 & Test 11 & \PASS{} (68~kg modeled) & ACCEPTED \\
}

\pagebreak

\section{\unclass Basis of Comparison}

\tableThrNow{\unclasslabel}{\unclass Basis of Comparison}{accrpt-boc}{l,2in}{l,1.5in}{l,2.5in}
{
    Source & Type & Parameters Validated \\
    AAIB Investigation Report [4] & Accident investigation & MTOW, thrust, configuration \\
    US Patent US20240262499A1 [5] & Patent filing & Battery configuration, voltage, capacity \\
    T-Motor U15II KV100 Specs [6] & Manufacturer data & Motor KV, max power, dimensions \\
    Cheeseman \& Bennett (1955) [8] & Published research & Ground effect thrust augmentation \\
    Johnson (1980) [9] & Textbook & Helicopter rotor aerodynamics, VRS \\
    Leishman (2006) [10] & Textbook & Rotorcraft aerodynamics, momentum theory \\
    TRV-150 Product Data [11] & Manufacturer marketing & Cruise speed, max payload, endurance \\
}

\pagebreak

\section{\unclass References}

\unclass [1] DoD Instruction 5000.61, ``DoD M\&S VV\&A,'' September 17, 2024.

\unclass [2] MIL-STD-3022, ``Documentation of VV\&A for Models and Simulations,'' January 28, 2008.

\unclass [3] Army Directive 2024-02: Software Development, 2024.

\unclass [4] AAIB, ``Investigation to Malloy Aeronautics T-150, South Scarle, Nottinghamshire,'' 2023.

\unclass [5] Malloy Aeronautics, ``Battery Pack Topology for UAV,'' US Patent US20240262499A1, 2024.

\unclass [6] T-Motor, ``U15II KV100 Motor Specifications,'' 2024.

\unclass [7] T-150 UAS AFSIM Plugin Suite, ``ADD v3.0,'' March 5, 2026.

\unclass [8] I.C. Cheeseman and W.E. Bennett, ``The Effect of Ground on a Helicopter Rotor,'' ARC R\&M 3021, 1955.

\unclass [9] W. Johnson, ``Helicopter Theory,'' Princeton University Press, 1980.

\unclass [10] J.G. Leishman, ``Principles of Helicopter Aerodynamics,'' 2nd ed., Cambridge, 2006.

\unclass [11] SURVICE / Malloy Aeronautics, ``TRV-150 Heavy-Lift UAS Product Data Sheet,'' 2024.

\unclass [12] NGIC, ``Impact of AI-Assisted Development on M\&S: T-150 UAS Case Study,'' March 4, 2026.

\pagebreak

\section{\unclass Acronyms}

\tableTwoNow{\unclasslabel}{\unclass Acronyms}{accrpt-acronyms}{l,1.5in}{l,4.5in}
{
    Acronym & Definition \\
    ADD & Architecture Description Document \\
    AFSIM & Advanced Framework for Simulation, Integration, and Modeling \\
    AGL & Above Ground Level \\
    BEMT & Blade Element Momentum Theory \\
    CG & Center of Gravity \\
    C-UAS & Counter-Unmanned Aerial System \\
    CONOPS & Concept of Operations \\
    DLL & Dynamic Link Library \\
    DoD & Department of Defense \\
    DoDI & DoD Instruction \\
    6DOF & Six Degrees of Freedom \\
    ESC & Electronic Speed Controller \\
    GE & Ground Effect \\
    ICD & Interface Control Document \\
    IGE & In Ground Effect \\
    ISA & International Standard Atmosphere \\
    KV & Motor velocity constant (RPM per Volt) \\
    M\&S & Modeling and Simulation \\
    MTOW & Maximum Takeoff Weight \\
    NaN & Not a Number \\
    NGIC & National Ground Intelligence Center \\
    OCV & Open-Circuit Voltage \\
    OGE & Out of Ground Effect \\
    OSINT & Open-Source Intelligence \\
    PID & Proportional-Integral-Derivative \\
    RPM & Revolutions Per Minute \\
    SME & Subject Matter Expert \\
    SOC & State of Charge \\
    TDD & Test Design Document \\
    UAM & User and Analyst Manual \\
    UAS & Unmanned Aerial System \\
    V\&V & Verification and Validation \\
    VRS & Vortex Ring State \\
    VV\&A & Verification, Validation, and Accreditation \\
}

\pagebreak

\section{\unclass Glossary}

\tableTwoNow{\unclasslabel}{\unclass Glossary}{accrpt-glossary}{l,1.5in}{l,4.5in}
{
    Term & Definition \\
    Accreditation & Official certification that a model or simulation is acceptable for use for a specific purpose. Per DoDI 5000.61. \\
    AuxData & AFSIM's platform-level key-value attribute store used for inter-processor communication. \\
    Blackboard architecture & A software pattern where components communicate by reading and writing to a shared data store. \\
    Calibration gap & A quantitative discrepancy between model output and reference data that is understood and documented. \\
    Conditional accreditation & An accreditation decision approving use subject to specified conditions and limitations. \\
    Coulomb counting & A battery SOC estimation method that integrates current over time. \\
    Ground effect & The increase in rotor thrust efficiency when operating near the ground. \\
    Momentum theory & An analytical framework for rotor performance based on conservation of momentum. \\
    Parametric-level fidelity & A level where key system parameters are represented with values traceable to source data. \\
    Vortex Ring State & A hazardous condition when a rotorcraft descends into its own downwash. \\
}

\pagebreak

\section{\unclass Accreditation Plan (Reference)}

\unclass The Accreditation Plan for the T-150 UAS AFSIM Simulation Model v2.0 is maintained as a separate document. Refer to the Accreditation Plan v1.0 for accreditation strategy, required evidence, assessment criteria, and decision framework.

\section{\unclass V\&V Report (Reference)}

\unclass The V\&V Report for the T-150 UAS AFSIM Simulation Model v2.0 is maintained as a separate document. The report contains detailed test procedures, raw results, pass/fail determinations, and analysis supporting the requirements assessment in Section 2.

\pagebreak

%% ============================================================
%%  SIGNATURE BLOCK
%% ============================================================

\section*{\unclass Signature Block}
\addcontentsline{toc}{section}{Signature Block}

\tableTwoNow{\unclasslabel}{\unclass Accreditation Authority Signature}{accrpt-sig-auth}{l,2.5in}{l,3.5in}
{
    Field & Value \\
    Accreditation Authority & [PLACEHOLDER - Name, Title, Organization] \\
    Signature & [PLACEHOLDER] \\
    Date & [PLACEHOLDER] \\
    Accreditation Decision & \textbf{CONDITIONAL ACCREDITATION} \\
    Effective Date & [PLACEHOLDER] \\
    Expiration / Review Date & [PLACEHOLDER - Typically 2 years from effective date] \\
}

\tableTwoNow{\unclasslabel}{\unclass V\&V Agent Signature}{accrpt-sig-vv}{l,2.5in}{l,3.5in}
{
    Field & Value \\
    V\&V Agent & [PLACEHOLDER - Name, Title, Organization] \\
    Signature & [PLACEHOLDER] \\
    Date & [PLACEHOLDER] \\
    V\&V Recommendation & Conditional Accreditation recommended based on V\&V findings \\
}

\tableTwoNow{\unclasslabel}{\unclass M\&S Proponent Signature}{accrpt-sig-prop}{l,2.5in}{l,3.5in}
{
    Field & Value \\
    M\&S Proponent & [PLACEHOLDER - Name, Title, Organization] \\
    Signature & [PLACEHOLDER] \\
    Date & [PLACEHOLDER] \\
    Concurrence & [PLACEHOLDER] \\
}

% glossary print function
\doPrintGloss

\doPrintBib%

\newpage\label{mylastpage}
\end{document}

%%% - - -   UNCLASSIFIED   - - - %%%
